%File: anonymous-submission-latex-2026.tex
\documentclass[letterpaper]{article} % DO NOT CHANGE THIS
\usepackage[submission]{aaai2026}  % DO NOT CHANGE THIS
\usepackage{times}  % DO NOT CHANGE THIS
\usepackage{helvet}  % DO NOT CHANGE THIS
\usepackage{courier}  % DO NOT CHANGE THIS
\usepackage[hyphens]{url}  % DO NOT CHANGE THIS
\usepackage{graphicx} % DO NOT CHANGE THIS
\urlstyle{rm} % DO NOT CHANGE THIS
\def\UrlFont{\rm}  % DO NOT CHANGE THIS
\usepackage{natbib}  % DO NOT CHANGE THIS AND DO NOT ADD ANY OPTIONS TO IT
\usepackage{caption} % DO NOT CHANGE THIS AND DO NOT ADD ANY OPTIONS TO IT
\frenchspacing  % DO NOT CHANGE THIS
\setlength{\pdfpagewidth}{8.5in} % DO NOT CHANGE THIS
\setlength{\pdfpageheight}{11in} % DO NOT CHANGE THIS
%
% These are recommended to typeset algorithms but not required. See the subsubsection on algorithms. Remove them if you don't have algorithms in your paper.
\usepackage{algorithm}
\usepackage{algorithmic}

\usepackage{amsmath} 
\usepackage{amsthm}
\newtheorem{definition}{Definition}  
\newtheorem{theorem}{Theorem}
\newtheorem{lemma}{Lemma}
\usepackage{booktabs}
\usepackage{multirow}

%
% These are are recommended to typeset listings but not required. See the subsubsection on listing. Remove this block if you don't have listings in your paper.
\usepackage{newfloat}
\usepackage{listings}
\DeclareCaptionStyle{ruled}{labelfont=normalfont,labelsep=colon,strut=off} % DO NOT CHANGE THIS
\lstset{%
	basicstyle={\footnotesize\ttfamily},% footnotesize acceptable for monospace
	numbers=left,numberstyle=\footnotesize,xleftmargin=2em,% show line numbers, remove this entire line if you don't want the numbers.
	aboveskip=0pt,belowskip=0pt,%
	showstringspaces=false,tabsize=2,breaklines=true}
\floatstyle{ruled}
\newfloat{listing}{tb}{lst}{}
\floatname{listing}{Listing}
%
% Keep the \pdfinfo as shown here. There's no need
% for you to add the /Title and /Author tags.
\pdfinfo{
/TemplateVersion (2026.1)
}

\usepackage{amsfonts}

% DISALLOWED PACKAGES
% \usepackage{authblk} -- This package is specifically forbidden
% \usepackage{balance} -- This package is specifically forbidden
% \usepackage{color (if used in text)
% \usepackage{CJK} -- This package is specifically forbidden
% \usepackage{float} -- This package is specifically forbidden
% \usepackage{flushend} -- This package is specifically forbidden
% \usepackage{fontenc} -- This package is specifically forbidden
% \usepackage{fullpage} -- This package is specifically forbidden
% \usepackage{geometry} -- This package is specifically forbidden
% \usepackage{grffile} -- This package is specifically forbidden
% \usepackage{hyperref} -- This package is specifically forbidden
% \usepackage{navigator} -- This package is specifically forbidden
% (or any other package that embeds links such as navigator or hyperref)
% \indentfirst} -- This package is specifically forbidden
% \layout} -- This package is specifically forbidden
% \multicol} -- This package is specifically forbidden
% \nameref} -- This package is specifically forbidden
% \usepackage{savetrees} -- This package is specifically forbidden
% \usepackage{setspace} -- This package is specifically forbidden
% \usepackage{stfloats} -- This package is specifically forbidden
% \usepackage{tabu} -- This package is specifically forbidden
% \usepackage{titlesec} -- This package is specifically forbidden
% \usepackage{tocbibind} -- This package is specifically forbidden
% \usepackage{ulem} -- This package is specifically forbidden
% \usepackage{wrapfig} -- This package is specifically forbidden
% DISALLOWED COMMANDS
% \nocopyright -- Your paper will not be published if you use this command
% \addtolength -- This command may not be used
% \balance -- This command may not be used
% \baselinestretch -- Your paper will not be published if you use this command
% \clearpage -- No page breaks of any kind may be used for the final version of your paper
% \columnsep -- This command may not be used
% \newpage -- No page breaks of any kind may be used for the final version of your paper
% \pagebreak -- No page breaks of any kind may be used for the final version of your paperr
% \pagestyle -- This command may not be used
% \tiny -- This is not an acceptable font size.
% \vspace{- -- No negative value may be used in proximity of a caption, figure, table, section, subsection, subsubsection, or reference
% \vskip{- -- No negative value may be used to alter spacing above or below a caption, figure, table, section, subsection, subsubsection, or reference

\setcounter{secnumdepth}{0} %May be changed to 1 or 2 if section numbers are desired.

% The file aaai2026.sty is the style file for AAAI Press
% proceedings, working notes, and technical reports.
%

% Title

% Your title must be in mixed case, not sentence case.
% That means all verbs (including short verbs like be, is, using,and go),
% nouns, adverbs, adjectives should be capitalized, including both words in hyphenated terms, while
% articles, conjunctions, and prepositions are lower case unless they
% directly follow a colon or long dash
\title{Flatness Matters: Theory and Practice of Flat Minima in Multi-LoRA Continual Learning}
\author{
    %Authors
    % All authors must be in the same font size and format.
    Written by AAAI Press Staff\textsuperscript{\rm 1}\thanks{With help from the AAAI Publications Committee.}\\
    AAAI Style Contributions by Pater Patel Schneider,
    Sunil Issar,\\
    J. Scott Penberthy,
    George Ferguson,
    Hans Guesgen,
    Francisco Cruz\equalcontrib,
    Marc Pujol-Gonzalez\equalcontrib
}
\affiliations{
    %Afiliations
    \textsuperscript{\rm 1}Association for the Advancement of Artificial Intelligence\\
    % If you have multiple authors and multiple affiliations
    % use superscripts in text and roman font to identify them.
    % For example,

    % Sunil Issar\textsuperscript{\rm 2},
    % J. Scott Penberthy\textsuperscript{\rm 3},
    % George Ferguson\textsuperscript{\rm 4},
    % Hans Guesgen\textsuperscript{\rm 5}
    % Note that the comma should be placed after the superscript

    1101 Pennsylvania Ave, NW Suite 300\\
    Washington, DC 20004 USA\\
    % email address must be in roman text type, not monospace or sans serif
    proceedings-questions@aaai.org
%
% See more examples next
}

%Example, Single Author, ->> remove \iffalse,\fi and place them surrounding AAAI title to use it
\iffalse
\title{My Publication Title --- Single Author}
\author {
    Author Name
}
\affiliations{
    Affiliation\\
    Affiliation Line 2\\
    name@example.com
}
\fi

\iffalse
%Example, Multiple Authors, ->> remove \iffalse,\fi and place them surrounding AAAI title to use it
\title{My Publication Title --- Multiple Authors}
\author {
    % Authors
    First Author Name\textsuperscript{\rm 1},
    Second Author Name\textsuperscript{\rm 2},
    Third Author Name\textsuperscript{\rm 1}
}
\affiliations {
    % Affiliations
    \textsuperscript{\rm 1}Affiliation 1\\
    \textsuperscript{\rm 2}Affiliation 2\\
    firstAuthor@affiliation1.com, secondAuthor@affilation2.com, thirdAuthor@affiliation1.com
}
\fi


% REMOVE THIS: bibentry
% This is only needed to show inline citations in the guidelines document. You should not need it and can safely delete it.
% \usepackage{bibentry}
% END REMOVE bibentry

\begin{document}

\maketitle

\begin{abstract}
AAAI creates proceedings, working notes, and technical reports directly from electronic source furnished by the authors. To ensure that all papers in the publication have a uniform appearance, authors must adhere to the following instructions.
\end{abstract}

% Uncomment the following to link to your code, datasets, an extended version or similar.
% You must keep this block between (not within) the abstract and the main body of the paper.
% \begin{links}
%     \link{Code}{https://aaai.org/example/code}
%     \link{Datasets}{https://aaai.org/example/datasets}
%     \link{Extended version}{https://aaai.org/example/extended-version}
% \end{links}


\section{Flatness Theory: A PAC-Bayes Perspective}



\bibliography{aaai2026}


\appendix

\section{Appendix A: Proof of the theorem}

\textbf{Theorem: Multi-Task PAC-Bayes Generalization Bound for Incremental LoRA with Chain-Dependent Hyperdistributions}


Consider a continual learning scenario with $n$ sequential tasks under the Incremental LoRA framework, subject to the following conditions:

1. Orthogonality of Parameter Spaces: For each task $t$, the LoRA parameter block $\Delta\mathbf{W}_t$ is constrained to be orthogonal to all previous parameter blocks, i.e., for any $i < t$, $A_i^\top A_t = 0$, where $A_i, A_t$ denote the adapter matrices of tasks $i$ and $t$, respectively.

2. Gaussian Posterior Assumption: The posterior parameter distribution for task $t$ is Gaussian, $Q_t = \mathcal{N}(\hat{\boldsymbol{\theta}}_t, \rho_t^2 \mathbf{I})$, where $\hat{\boldsymbol{\theta}}_t$ is the optimal parameter for task $t$, and $\rho_t^2$ reflects the posterior uncertainty.

3. Chain-Structured Hyperdistribution: The hyperprior for task $t$, $\mathcal{P}_t$, is conditioned on the cumulative hyperposterior of all previous tasks, i.e., the joint hyperprior is
   $$
   \mathcal{P}_{1:n} = \mathcal{P}_1 \otimes \mathcal{P}_2(\cdot|\mathcal{Q}_1) \otimes \cdots \otimes \mathcal{P}_n(\cdot|\mathcal{Q}_{1:n-1}),
   $$ where $\mathcal{Q}_{1:t-1}$ denotes the joint hyperposterior for tasks $1$ through $t-1$.

4. Sharpness Definition: The expected sharpness for task $t$ is defined as
\begin{align}
&\mathrm{Sharpness}_t(\mathbf{W}_{t-1} + \hat{\Delta\mathbf{W}}_t,\, \rho_t^2) \notag\\
& = \mathbb{E}_{\epsilon \sim \mathcal{N}(0,\, \rho_t^2 \mathbf{I})}
\left[ \widehat{L}_{S_t}(\mathbf{W}_{t-1} + \hat{\Delta\mathbf{W}}_t + \epsilon) \right] \notag\\
&\quad - \widehat{L}_{S_t}(\mathbf{W}_{t-1} + \hat{\Delta\mathbf{W}}_t) \notag
\end{align}
quantifying the increase in empirical loss due to parameter perturbation.

\begin{theorem}[PAC-Bayes Generalization Bound for Incremental Multi-Task LoRA]
With probability at least $1-\delta$, the generalization risk $L(\mathcal{Q}*{1:n})$ for $n$ sequential tasks under the incremental LoRA framework satisfies
\begin{align}
L(\mathcal{Q}_{1:n}) \leq\; & \frac{1}{n} \sum_{t=1}^n 
    \mathbb{E}_{P_{1:n} \sim \mathcal{Q}_{1:n}}
    \Big[
        \widehat{L}_{S_t}(\mathbf{W}_{t-1} + \hat{\Delta\mathbf{W}}_t)
    \notag \\
&\qquad +\; \mathrm{Sharpness}_t(\mathbf{W}_{t-1} + \hat{\Delta\mathbf{W}}_t,\, \rho_t^2)
    \Big] 
    \notag \\
&\quad + (b-a) \left\{ 
    \frac{KL_n^{\text{task}} + KL_n^{\text{hyper}} + \log(m/\delta)}
         {2(\bar{m}_n - 1)}
\right\}^{1/2}
\label{eq:lora-pac-bayes-bound}
\end{align}
where:
\begin{align}
&KL_n^{\mathrm{task}} = \sum_{t=1}^n KL(Q_t\|P_t), \notag  \\
&KL_n^{\mathrm{hyper}} = KL(\mathcal{Q}_1 \| \mathcal{P}_1) + \sum_{t=2}^n \mathbb{E}_{\mathcal{Q}_{1:t-1}}\big[KL(\mathcal{Q}_t \| \mathcal{P}_t(\cdot \| \mathcal{Q}_{1:t-1}))\big], \notag\\
&\bar{m}_n = n \big/ \left(\sum_{t=1}^n 1/m_t\right), \notag
\end{align}
with $m_t$ denoting the sample size for task $t$, and $b-a$ the range of the bounded loss function.
\end{theorem}

\begin{proof}
\textbf{Step 1: Single-task PAC-Bayes Bound (0→1)}

The PAC-Bayesian theory provides a theoretical foundation for analyzing the generalization properties of stochastic model selection, where algorithms probabilistically select models based on a probability distribution over the hypothesis space  $\mathcal{F} $. 


\begin{theorem} \label{the: PAC1999}
\citep{mcallesterPACBayesianStochasticModel2003}
For loss functions $\ell^{'}$ mapping to the range $[0,1]$, the following generalization bound holds for any prior $P$ over parameters with probability $1 - \delta$ over the choice of the training set $S$, for any posterior $Q$ over parameters.
\begin{equation}\label{eq: PAC2003}
     \forall^{\delta}S,\ \forall Q,\ \mathbb{E}_{f \in Q}[L_{\mathcal{D}}^{\ell^{'}}(f] \leq  \mathbb{E}_{f \in Q}[\hat{L}^{\ell^{'}}_{S}(f)] +\sqrt{\frac{KL(Q\|P)+\log \frac{m}{\delta}}{2(m-1)}}
\end{equation}
where $m=|S|$ is the number of sample in the training set $S$ and $f$ denotes a prediction function sampled from the posterior distribution $Q$ over the hypothesis space $\mathcal{F}$.
\end{theorem}

Theorem \ref{the: PAC1999} is limited to loss functions $\ell^{'}$ mapping to the range $[0,1]$. According to \citep{germainPACBayesianTheoryMeets2016}, by straightforward rescaling, we can extend it to any bounded loss, i.e., $\ell: \mathcal{F} \times \mathcal{X} \times \mathcal{Y} \rightarrow[a, b]$, where $[a, b] \subset \mathbb{R}$. This is done using $\beta:=b-a$ and the rescaled loss function $\ell^{\prime}(f, x, y):=(\ell(f, x, y)-a) /(b-a) \in[0,1]$. After few arithmetic manipulations, we can rewrite Equation (\ref{eq: PAC2003}) as
\begin{equation}
\begin{aligned}
        &\mathbb{E}_{f \sim Q}[{L}_\mathcal{D}^{\ell}(f)] \leq \mathbb{E}_{f\sim Q}[\widehat{{L}}_S^{\ell}(f)] \\
        &+ (b - a) \cdot \sqrt{\frac{\mathrm{KL}(Q \| P) + \log \frac{m}{\delta}}{2(m - 1)}} .
\end{aligned}
\end{equation}
If the right-hand side of the bound exceeds $b-a$, the bound reduces to
$\mathbb{E}_{f \sim Q}[L_D^{\ell}(f)] \leq b,$
which is always true by definition. So we focus on the nontrivial regime where the complexity term is less than $b-a$, ,which equals to
$\sqrt{\frac{KL(Q\|P)+\log \frac{m}{\delta}}{2(m-1)}} <1. $

% and thus offers no additional information about model's generalization ability. Therefore, only when the upper bound is strictly less than $b$ does it serve as a nontrivial guarantee about the model's generalization performance. So we should only care when the complexity term no more than $b-a$,which equals to
% $\sqrt{\frac{KL(Q\|P)+\log \frac{m}{\delta}}{2(m-1)}} <1. $

% Following \citep{pentinaPACBayesianBoundLifelong2014}, we introduce hyperdistributions $\mathcal{P}$ (hyperprior) and $\mathcal{Q}$ (hyperposterior), where each parameter prior $P$ is a sample from $\mathcal{P}$, and $\mathcal{Q}$ represents the posterior over parameter priors after observing data. In PAC-Bayes theory with hyperdistributions, $Q_\text{joint}(f, P) = Q(P)\mathcal{Q}(P)$ represents the learned joint distribution over parameters and hyperparameters, where $\mathcal{Q}(P)$ is the hyperposterior (after observing the data), and $Q(P)$ is the parameter posterior conditioned on $P$. The reference joint distribution $P_\text{joint}(f, P) = P(f)\mathcal{P}(P)$ encodes the prior belief before observing the data, where $\mathcal{P}(P)$ is the hyperprior and $P(f)$ is the parameter prior under hyperparameter $P$. The KL divergence between these two joint distributions quantifies both the adaptation to data in parameter space and in hyperparameter space.  $KL(Q_\text{joint} \| P_\text{joint}) = \mathbb{E}_{P \sim \mathcal{Q}} \left[ KL(Q(P) \| P) \right] + KL(\mathcal{Q} \| \mathcal{P}).$

Following \citep{pentinaPACBayesianBoundLifelong2014}, we introduce the notion of hyperdistributions. Let $\mathcal{P}$ denote the hyperprior (a distribution over parameter priors $P$) and $\mathcal{Q}$ denote the hyperposterior, the updated distribution over priors after observing data. For each choice of prior $P$, the parameter posterior is denoted $Q(f|P)$, while $P(f)$ is the corresponding parameter prior under $P$.

In PAC-Bayes theory with hyperdistributions, the learned joint distribution over parameters and hyperparameters is given by $Q_\text{joint}(f, P) = Q(f|P)\mathcal{Q}(P)$, and the reference joint distribution before seeing data is $P_\text{joint}(f, P) = P(f)\mathcal{P}(P)$. The KL divergence between these two joint distributions quantifies the adaptation to data in both parameter space and hyperparameter space, and can be decomposed as:
$$KL(Q_\text{joint} \| P_\text{joint}) = \mathbb{E}_{P \sim \mathcal{Q}} \left[ KL(Q(f|P) \| P(f)) \right] + KL(\mathcal{Q} \| \mathcal{P}).$$


Taking the expectation over the hyperposterior $P \sim \mathcal{Q}$, the PAC-Bayes bound for hyperdistributions is:
\begin{equation}
\begin{aligned}
        &\mathbb{E}_{P \sim \mathcal{Q}}\left[\, \mathbb{E}_{f \sim Q(P)} [L_{\mathcal{D}}(f)] \,\right] 
        \leq \mathbb{E}_{P \sim \mathcal{Q}}\left[\, \mathbb{E}_{f \sim Q(P)} [\hat{L}_{S}(f)] \,\right]  \\
        & + (b-a)\sqrt{ \frac{ \mathbb{E}_{P \sim \mathcal{Q}}\left[ KL(Q(P) \| P) \right] + KL(\mathcal{Q} \| \mathcal{P}) + \log \frac{m}{\delta} }{2(m-1)} }
\end{aligned}
\end{equation}


Now, let $Q = \mathcal{N}(\hat{\boldsymbol{\theta}}, \rho^2 I)$, where $\hat{\boldsymbol{\theta}}$ is the trained parameter vector. Then,
$$\mathbb{E}_{\boldsymbol{\theta} \sim Q}[\hat{L}_S(\boldsymbol{\theta})] = \mathbb{E}_{\boldsymbol{\epsilon} \sim \mathcal{N}(0, \rho^2 I)}[\hat{L}_S(\hat{\boldsymbol{\theta}} + \boldsymbol{\epsilon})].$$
Refer the \textbf{expected sharpness} definition (\ref{def: Task-Aware Expected Sharpness}):
$$\mathrm{Sh}_S(\hat{\boldsymbol{\theta}}, \rho^2) := \mathbb{E}_{\boldsymbol{\epsilon} \sim \mathcal{N}(0, \rho^2 I)}\left[ \hat{L}_S(\hat{\boldsymbol{\theta}} + \boldsymbol{\epsilon}) \right] - \hat{L}_S(\hat{\boldsymbol{\theta}})$$
Thus,
$$\mathbb{E}_{\boldsymbol{\theta} \sim Q}[\hat{L}_S(\boldsymbol{\theta})] = \hat{L}_S(\hat{\boldsymbol{\theta}}) + \mathrm{Sh}_S(\hat{\boldsymbol{\theta}}, \rho^2)$$

In the LoRA framework, the full-model parameters are $\boldsymbol{W}_1 = \boldsymbol{W}_0 + \Delta\boldsymbol{W}_1$, where $\boldsymbol{W}_0$ denotes the frozen pre-trained backbone and $\Delta\boldsymbol{W}_1$ is the task-specific low-rank adaptation. The full-model distributions factorize as $Q_1^{\text{full}} = \delta_{\boldsymbol{W}_0} \otimes Q_1$ and $P_1^{\text{full}} = \delta_{\boldsymbol{W}_0} \otimes P_1$, leading to:
\begin{equation}
    \text{KL}^{\text{task}}(Q_1^{\text{full}} \| P_1^{\text{full}}) = \text{KL}^{\text{task}}(Q_1 \| P_1),
\end{equation}
where $Q_1$ and $P_1$ are distributions over $\Delta\boldsymbol{W}_1$. 


Although the posterior $Q$ is defined only over $\Delta\boldsymbol{W}_1$, the model prediction and corresponding loss are determined by the entire effective parameter $\boldsymbol{W}_0 + \Delta\boldsymbol{W}_1$. Therefore, when evaluating sharpness, the perturbation is applied to the full parameter vector $\boldsymbol{W}_0 + \Delta\boldsymbol{W}_1$, not just to the low-rank increment $\Delta\boldsymbol{W}_1$. Formally, the expected sharpness in the learned model $\boldsymbol{W}_0 + \hat{\Delta\boldsymbol{W}}_1$ is defined as:
\begin{equation}
\begin{aligned}
&\mathrm{Sh}_S\big(\boldsymbol{W}_0 + \hat{\Delta\mathbf{W}}_1, \rho^2\big) := \\
&\mathbb{E}_{\boldsymbol{\epsilon} \sim \mathcal{N}(0, \rho^2 I)} \Big[
    \widehat{L}_S\big(\boldsymbol{W}_0 + \hat{\Delta\mathbf{W}}_1 + \boldsymbol{\epsilon}\big)
\Big] - \widehat{L}_S\big(\boldsymbol{W}_0 + \hat{\Delta\mathbf{W}}_1\big) 
\notag
\end{aligned}
\end{equation}
where the perturbation $\boldsymbol{\epsilon}$ is added to the full parameter vector. Thus,
\begin{equation}
    \begin{aligned}
    &\mathbb{E}_{\Delta\mathbf{W} \sim Q}[\widehat{L}_S(\mathbf{W}_0 + \Delta\mathbf{W}_1)] \\
    &= \mathbb{E}_{\boldsymbol{\epsilon} \sim \mathcal{N}(0, \rho^2 I)}[\widehat{L}_S(\mathbf{W}_0 + \hat{\Delta\mathbf{W}}_1 + \boldsymbol{\epsilon})]  \\
    &=\widehat{L}_S(\mathbf{W}_0 + \hat{\Delta\mathbf{W}}_1) + \mathrm{Sh}_S(\mathbf{W}_0+\hat{\Delta\mathbf{W}}_1, \rho^2)
\end{aligned}
\end{equation}


Finally, the corresponding PAC-Bayes bound, incorporating sharpness under the LoRA setting, is as follows:

\begin{align}
& \mathbb{E}_{P \sim \mathcal{Q}}\, \mathbb{E}_{\Delta\mathbf{W} \sim Q(P)} 
    \left[ L_{\mathcal{D}}\left(\boldsymbol{W}_0 + \Delta\mathbf{W}_1\right) \right] \notag \\
& \leq\; \mathbb{E}_{P \sim \mathcal{Q}} \Big[
        \widehat{L}_S\big(\boldsymbol{W}_0 + \hat{\Delta\mathbf{W}}_1\big)
        + \mathrm{Sh}_S\big(\boldsymbol{W}_0 + \hat{\Delta\mathbf{W}}_1,\, \rho^2\big)
    \Big] \notag \\
& + (b-a) 
    \Bigg\{ \frac{1}{2(m-1)} \bigg[
        KL^{\text{task}}_1 + KL^{\text{hyper}}_1 + \log\frac{m}{\delta}
    \bigg] \Bigg\}^{1/2}
\label{eq:single-task-pac-bayes}
\end{align}
where $KL^{\text{task}} = \mathbb{E}_{P \sim \mathcal{Q}}\left[ KL(Q(P)|P) \right]$,
$KL^{\text{hyper}} = KL(\mathcal{Q}|\mathcal{P})$,
and $\log(m/\delta)$ is the parameter term.


\textbf{Step 2: Extending the PAC-Bayes Bound from One Task to Two Tasks  (1→2)}
To extend the result to two sequential tasks ($n=2$), we consider the incremetnal LoRA scenario in which the LoRA parameter block for the second task, $\Delta\mathbf{W}_2$, is trained independently, with the backbone fixed at $\mathbf{W}_1 = \mathbf{W}_0 + \hat{\Delta\mathbf{W}}_1$. 


In incremental LoRA, we consider the parameter increments for each task, ${\Delta\mathbf{W}}_{t=1}^n$, as forming mutually orthogonal subspaces in the parameter space. Formally, let $U_t$ denote the subspace spanned by the LoRA parameter block for task $t$. The orthogonality condition requires that for any pair of tasks $i \neq t$,
$\langle u, w \rangle = 0, \quad \forall\, u \in U_i,\, w \in U_t,$
where $\langle \cdot, \cdot \rangle$ denotes the inner product in the parameter space. Therefore, achieving orthogonality between the LoRA subspaces of task $i$ and task $t$ can be expressed as
$${O}_{i,t} = A_i^{\top} A_t = 0,$$
where $A_i$ and $A_t$ are the adaptation matrices for tasks $i$ and $t$, respectively. Under this assumption, the joint KL divergence of product measures satisfies:   \begin{equation}
      KL\left(\bigotimes_{t=1}^n Q_t \, \bigg\| \, \bigotimes_{t=1}^n P_t\right) = \sum_{t=1}^n KL(Q_t \| P_t).
  \end{equation}
where $Q_t$ and $P_t$ denote the posterior and prior distributions over the LoRA parameters for task $t$.



The corresponding joint hyperprior and hyperposterior for the two-task setting are denoted as $\mathcal{P}_{1:2}$ and $\mathcal{Q}_{1:2}$. Here, $\mathcal{P}_{1:2}$ represents the joint hyperprior over both tasks, where the hyperprior for the second task is conditionally dependent on the hyperposterior of the first task, i.e.,
$\mathcal{P}_{1:2} = \mathcal{P}_1 \otimes \mathcal{P}_2(\cdot\,|\,\mathcal{Q}_1).$
Similarly, the joint hyperposterior is given by
$\mathcal{Q}_{1:2} = \mathcal{Q}_1 \otimes \mathcal{Q}_2,$
assuming conditional independence between the task-wise hyperposteriors. Because $\mathcal{P}_2$ is conditionally dependent on $\mathcal{Q}_1$, the joint KL divergence between the hyperdistributions admits a chain-structured decomposition:
\begin{equation}
\begin{aligned}
    KL^{\text{hyper}}_2&=KL(\mathcal{Q}_{1:2}\,\|\,\mathcal{P}_{1:2})  \\
    &= KL(\mathcal{Q}_1\,\|\,\mathcal{P}_1) + \mathbb{E}_{\mathcal{Q}_1}\left[ KL(\mathcal{Q}_2\,\|\,\mathcal{P}_2(\cdot\,|\,\mathcal{Q}_1)) \right].
\end{aligned}
\end{equation}
This structure reflects that the uncertainty or inductive bias for the second task is informed by the knowledge gained from the first task. In practical terms, the conditional hyperprior $\mathcal{P}_2(\cdot,|,\mathcal{Q}_1)$ can take the form of a Gaussian distribution whose parameters (e.g., covariance) are derived from the posterior statistics of the previous task. For example,
$P_2 = \mathcal{N}\big(\mathbf{0},\, \mathbf{F}_1^{-1}\big),$
where $\mathbf{F}_1$ is the Fisher information matrix estimated from the first task, thus encoding the structure and uncertainty learned in the previous stage.


Analogously, for the parameter-level distributions, the KL divergence decomposes as
\begin{equation}
    KL^{\text{task}}_2=KL(Q_{1:2}\,\|\,P_{1:2}) = KL(Q_1\,\|\,P_1) + KL(Q_2\,\|\,P_2),
\end{equation}
where $P_2$ is sampled from the conditional hyperprior $\mathcal{P}_2(\cdot|\mathcal{Q}_1)$.



Therefore, the generalization bound for two sequential tasks with recursive, dependent hyperdistributions takes the following form:

\begin{align}
L(\mathcal{Q}_{1:2}) \leq\; & \frac{1}{2} \sum_{t=1}^2 
    \mathbb{E}_{P_{1:2} \sim \mathcal{Q}_{1:2}} 
    \Big[\, 
        \widehat{L}_{S_t}(\mathbf{W}_{t-1} + \hat{\Delta\mathbf{W}}_t)
    \notag \\
& \qquad 
    +\; \mathrm{Sh}_t\big(\mathbf{W}_{t-1} + \hat{\Delta\mathbf{W}}_t,\, \rho_t^2\big)
    \Big] \notag \\
& + (b-a)\, 
    \Bigg\{ \frac{1}{2(\bar{m}_2 - 1)} \bigg[
        KL^{\text{task}}_2 + KL^{\text{hyper}}_2 + \log\frac{m}{\delta}
    \bigg] \Bigg\}^{1/2}
\label{eq:cl-bound-chain-compact}
\end{align}

where 

\begin{align*}
&L(\mathcal{Q}_{1:2}) :=\;  \mathbb{E}_{P_{1:2} \sim \mathcal{Q}_{1:2}}
\bigg\{
    \frac{1}{2} \Big[
        \mathbb{E}_{\Delta\mathbf{W}_1 \sim Q_1(P_1)}\big[
            L_{\mathcal{D}_1}(\mathbf{W}_0 + \Delta\mathbf{W}_1)
        \big] \notag \\
& \qquad
        +
        \mathbb{E}_{\Delta\mathbf{W}_2 \sim Q_2(P_2)}\big[
            L_{\mathcal{D}_2}(\mathbf{W}_1 + \Delta\mathbf{W}_2)
        \big]
    \Big]
\bigg\} \\
&KL^{\text{task}}_2 = 
    \mathbb{E}_{P_1 \sim \mathcal{Q}_1}\left[ KL(Q_1(P_1)\|P_1) \right]
    + \mathbb{E}_{P_2 \sim \mathcal{Q}_2}\left[ KL(Q_2(P_2)\|P_2) \right] \\
&KL^{\text{hyper}}_2 = 
    KL(\mathcal{Q}_1\|\mathcal{P}_1)
    + \mathbb{E}_{\mathcal{Q}_1}\left[ KL(\mathcal{Q}_2 \| \mathcal{P}_2(\cdot\,|\,\mathcal{Q}_1)) \right] \\
&\bar{m}_2 = 2 / \left(\frac{1}{m_1} + \frac{1}{m_2}\right)
\end{align*}


 All terms involving empirical risk, sharpness, and complexity are averaged over the joint hyperposterior $\mathcal{Q}{1:2}$. The chain-structured KL terms reflect that each new task's hyperprior may depend on the hyperposterior of previous tasks, capturing the transfer and adaptation of uncertainty and prior structure in incremental learning.


\textbf{Step 3: Multi-task PAC-Bayes Bound (n→n+1)}

Assume the following holds for $n$ sequential tasks:


\textbf{Parameter Distribution KL:}
    \begin{align}
     KL^{\text{task}}_n=KL(Q_{1:n} \| P_{1:n}) 
    &= \sum_{t=1}^n KL(Q_t \| P_t)
    \end{align}

 
\textbf{Hyperdistribution KL:}
    \begin{align}
    &KL^{\text{hyper}}_n=KL(\mathcal{Q}_{1:n} \| \mathcal{P}_{1:n})  \\
    &= KL(\mathcal{Q}_1 \| \mathcal{P}_1) \notag 
      + \sum_{t=2}^n 
        \mathbb{E}_{\mathcal{Q}_{1:t-1}} \left[ 
            KL\left(\mathcal{Q}_t \| \mathcal{P}_t\big(\cdot\,|\,\mathcal{Q}_{1:t-1}\big)\right) 
        \right]
    \end{align}

\textbf{Generalization Bound:}
    \begin{align}
    L(\mathcal{Q}_{1:n}) \leq\; & \frac{1}{n} \sum_{t=1}^n 
        \mathbb{E} \Big[
            \widehat{L}_{S_t} + \mathrm{Sh}_t
        \Big] \notag \\
    & + (b-a) \sqrt{
        \frac{KL^{\text{task}}_n +KL^{\text{hyper}}_n + \log(m/\delta)}{2(\bar{m}_n - 1)}
    }
    \end{align}
    where
    % \begin{align} 
    %   L(\mathcal{Q}_{1:n}) &:= \mathbb{E}_{P_{1:n} \sim \mathcal{Q}_{1:n}} \left\{ \frac{1}{n} \sum_{t=1}^n \mathbb{E}_{\Delta\mathbf{W}_t \sim Q_t(P_t)} \left[ L_{\mathcal{D}_t}(\mathbf{W}_{t-1} + \Delta\mathbf{W}_t) \right] \right\}   \notag  \\
    %   \bar{m}_n &= n \Big/ \left(\sum_{t=1}^n 1/m_t\right) \notag  
    % \end{align}
    \begin{align}
L(\mathcal{Q}_{1:n}) :=\; 
    &\frac{1}{n} \sum_{t=1}^n \;
      \mathbb{E}_{P_{1:n} \sim \mathcal{Q}_{1:n}} \;
      \mathbb{E}_{\Delta\mathbf{W}_t \sim Q_t(P_t)} 
      \big[
        L_{\mathcal{D}_t}(\mathbf{W}_{t-1} + \Delta\mathbf{W}_t)
      \big]  \notag \\
\bar{m}_n =\; 
    &n \bigg/ \Big( \sum_{t=1}^n \frac{1}{m_t} \Big) \notag
\end{align}


\textbf{Inductive Step: From $n$ to $n+1$ Tasks}


\textbf{(a) Extension of Orthogonality.}
The LoRA parameter block $\Delta\mathbf{W}_{n+1}$ is constrained to be orthogonal to all previous blocks:
    \begin{align}
        \forall\, i \leq n, \quad A_i^\top A_{n+1} = 0 \notag
    \end{align}
    Thus,
    \begin{equation}
        KL^{\text{task}}_{n+1}=KL(Q_{1:n+1} \| P_{1:n+1}) = \sum_{t=1}^{n+1} KL(Q_t \| P_t)
    \end{equation}

\textbf{(b) Chain-Structured Hyperprior.}
    The hyperprior for task $n+1$ is conditioned on the cumulative hyperposterior:
    \begin{align}
        \mathcal{P}_{1:n+1} = \mathcal{P}_{1:n} \otimes \mathcal{P}_{n+1}(\cdot\,|\,\mathcal{Q}_{1:n}) \notag
    \end{align}
    So,
    \begin{equation}
        \begin{aligned}
        KL^{\text{hyper}}_{n+1}
        &=KL(\mathcal{Q}_{1:n+1} \| \mathcal{P}_{1:n+1}) \\
        &= KL(\mathcal{Q}_{1:n} \| \mathcal{P}_{1:n}) +\, \mathbb{E}_{\mathcal{Q}_{1:n}}\left[ KL\big(\mathcal{Q}_{n+1} \| \mathcal{P}_{n+1}(\cdot\,|\,\mathcal{Q}_{1:n})\big) \right] 
        \\
        &= KL(\mathcal{Q}_1 \| \mathcal{P}_1) +
         \sum_{t=2}^{n+1} \mathbb{E}_{\mathcal{Q}_{1:t-1}}\left[ KL\left(\mathcal{Q}_t \| \mathcal{P}_t(\cdot\,|\,\mathcal{Q}_{1:t-1})\right) \right]
    \end{aligned}
    \end{equation}


\textbf{(c) Generalization Bound.}
    The empirical risk and sharpness terms for the new task are averaged in:
    \begin{align}
        L(\mathcal{Q}_{1:n+1}) \leq\; & \frac{1}{n+1} \sum_{t=1}^{n+1}
        \mathbb{E}_{P_{1:n+1} \sim \mathcal{Q}_{1:n+1}}\big[
            \widehat{L}_{S_t}(\mathbf{W}_{t-1} + \hat{\Delta\mathbf{W}}_t) \notag\\
        &\qquad +\, \mathrm{Sharpness}_t(\mathbf{W}_{t-1} + \hat{\Delta\mathbf{W}}_t,\, \rho_t^2)
        \big] \notag\\
        & + (b-a)\sqrt{
            \frac{KL_{n+1}^{\text{task}} + KL_{n+1}^{\text{hyper}} \log(m/\delta)}{2(\bar{m}_{n+1} - 1)}
        }
    \end{align}
    where
    \begin{equation}
        \begin{aligned}
        KL_{n+1}^{\text{task}} &= \sum_{t=1}^{n+1} KL(Q_t\|P_t)\\
        KL_{n+1}^{\text{hyper}} = & KL(\mathcal{Q}_1\|\mathcal{P}_1)  +\sum_{t=2}^{n+1} \mathbb{E}_{\mathcal{Q}_{1:t-1}}
        \left[ KL(\mathcal{Q}_t \| \mathcal{P}_t(\cdot\,|\,\mathcal{Q}_{1:t-1})) \right] \\
        \bar{m}_{n+1} &= (n+1)/\big(\sum_{t=1}^{n+1} 1/m_t\big) \notag
    \end{aligned}
    \end{equation}

\textbf{Step4: Conclusion General Bound for $n$ Tasks}

By induction, the PAC-Bayesian generalization bound for incremental multi-task LoRA with orthogonal blocks and recursively dependent hyperpriors is:
\begin{align}
        L(\mathcal{Q}_{1:n}) \leq\; & \frac{1}{n} \sum_{t=1}^{n}
        \mathbb{E}_{P_{1:n} \sim \mathcal{Q}_{1:n}}\big[
            \widehat{L}_{S_t}(\mathbf{W}_{t-1} + \hat{\Delta\mathbf{W}}_t) \notag\\
        &\qquad +\, \mathrm{Sh}_t(\mathbf{W}_{t-1} + \hat{\Delta\mathbf{W}}_t,\, \rho_t^2)
        \big] \notag\\
        & + (b-a)\sqrt{
            \frac{KL_{n}^{\text{task}} + KL_{n}^{\text{hyper}} \log(m/\delta)}{2(\bar{m}_{n} - 1)}
        }
    \end{align}


 This result demonstrates that, under the incremental LoRA framework, orthogonalization of parameter spaces and sequential modeling of hyperpriors allow the single-task PAC-Bayes theory to be strictly extended to multi-task continual learning. This structure ensures that model complexity, empirical risk, and sharpness are balanced while enabling transfer and retention of knowledge across tasks.




\end{proof}


% Inductive Step: From $n$ to $n+1$ Tasks

% \textbf{(a) Extension of Orthogonality.}
%  The LoRA parameter block $\Delta\mathbf{W}_{n+1}$ is constrained to be orthogonal to all previous blocks:
%  \begin{equation}
%  \forall i \leq n, ;; A_i^\top A_{n+1} = 0
%  \end{equation}
%  Thus,
%  \begin{equation}
%  KL(Q_{1:n+1} ,|, P_{1:n+1}) = \sum_{t=1}^{n+1} KL(Q_t ,|, P_t)
%  \end{equation}

% \textbf{(b) Chain-Structured Hyperprior.}
%  The hyperprior for task $n+1$ is conditioned on the cumulative hyperposterior of previous tasks:
%  \begin{equation}
%  \mathcal{P}*{1:n+1} = \mathcal{P}*{1:n} \otimes \mathcal{P}*{n+1}(\cdot|\mathcal{Q}*{1:n})
%  \end{equation}
%  So,
%  \begin{align}
%  & KL(\mathcal{Q}*{1:n+1} ,|, \mathcal{P}*{1:n+1}) =
%  KL(\mathcal{Q}*{1:n} ,|, \mathcal{P}*{1:n}) + \mathbb{E}*{\mathcal{Q}*{1:n}}\Big[ KL(\mathcal{Q}*{n+1} ,|, \mathcal{P}*{n+1}(\cdot|\mathcal{Q}*{1:n})) \Big] \notag \
%  &= KL(\mathcal{Q}\*1 ,|, \mathcal{P}\*1)
%  \+ \sum\*{t=2}^{n+1} \mathbb{E}\*{\mathcal{Q}*{1:t-1}}\big[ KL(\mathcal{Q}_t | \mathcal{P}*t(\cdot|\mathcal{Q}*{1:t-1})) \big]
%  \end{align}

% \textbf{(c) Generalization Bound.}
%  The empirical risk and sharpness term for the new task are averaged in:
%  \begin{align}
%  & L(\mathcal{Q}*{1:n+1}) \leq \frac{1}{n+1} \sum*{t=1}^{n+1} \mathbb{E}*{P*{1:n+1} \sim \mathcal{Q}*{1:n+1}}\left[
%  \widehat{L}*{S_t}(\mathbf{W}*{t-1} + \hat{\Delta\mathbf{W}}\*t)
%  \+ \mathrm{Sharpness}\*t(\mathbf{W}\*{t-1} + \hat{\Delta\mathbf{W}}\*t, \rho_t^2)
%  \right] \notag\
%  &\qquad + (b-a)\sqrt{\frac{\mathcal{K}\*{n+1} + \log(m/\delta)}{2(\bar{m}\*{n+1}-1)}}
%  \end{align}
%  where
%  \begin{align}
%  \mathcal{K}*{n+1} = \sum_{t=1}^{n+1} KL(Q_t|P_t) + KL(\mathcal{Q}*1|\mathcal{P}\*1)
%  \+ \sum\*{t=2}^{n+1} \mathbb{E}*{\mathcal{Q}_{1:t-1}}\left[ KL(\mathcal{Q}_t | \mathcal{P}*t(\cdot|\mathcal{Q}*{1:t-1})) \right]
%  \end{align}
%  and $\bar{m}_{n+1} = (n+1)/\big(\sum_{t=1}^{n+1} 1/m_t\big)$ is the harmonic mean of the sample sizes for all tasks.



% \textbf{Step 4: Inductive Conclusion General Bound for $n$ Tasks}


% By mathematical induction, the PAC-Bayesian generalization bound for incremental multi-task LoRA with orthogonal parameter blocks and recursively dependent hyperpriors is:
%  \begin{align}
%  L(\mathcal{Q}*{1:n}) \leq \frac{1}{n} \sum*{t=1}^n \mathcal{E}*t
%  \+ (b-a)\sqrt{\frac{\mathcal{K}\*n + \log(m/\delta)}{2(\bar{m}\*n-1)}}
%  \end{align}
%  where
%  \begin{equation}
%  \mathcal{E}\*t = \mathbb{E}\*{P\*{1:n} \sim \mathcal{Q}\*{1:n}}\big[
%  \widehat{L}*{S_t}(\mathbf{W}_{t-1} + \hat{\Delta\mathbf{W}}_t)
%  \+ \mathrm{Sharpness}*t(\mathbf{W}*{t-1} + \hat{\Delta\mathbf{W}}_t, \rho_t^2)
%  \big]
%  \end{equation}
%  and $\mathcal{K}_n$ is as defined above.



% \begin{align}

% \label{eq:cl_exp_risk-new}
% \end{align}

% \begin{align}
% L(\mathcal{Q}_{1:2}) \leq\; & \frac{1}{2}\sum_{t=1}^2
%     \mathbb{E}_{P_{1:2} \sim \mathcal{Q}_{1:2}} \Big[
%         \widehat{L}_{S_t}(\mathbf{W}_{t-1} + \hat{\Delta\mathbf{W}}_t)
%     \notag \\
% & \qquad
%         +\; \mathrm{Sharpness}_t\left(\mathbf{W}_{t-1} + \hat{\Delta\mathbf{W}}_t, \rho_t^2\right)
%     \Big] \notag \\
% & + (b-a) \left\{
%     \frac{1}{2(\bar{m}_2 - 1)} \Big[
%         \mathbb{E}_{P_1 \sim \mathcal{Q}_1}\big[ KL(Q_1(P_1)\|P_1) \big] \notag \\
% & \qquad
%         +\; \mathbb{E}_{P_2 \sim \mathcal{Q}_2}\big[ KL(Q_2(P_2)\|P_2) \big] \notag \\
% & \qquad
%         +\; KL(\mathcal{Q}_1\|\mathcal{P}_1)
%         +\; \mathbb{E}_{\mathcal{Q}_1}\left[ KL(\mathcal{Q}_2 \| \mathcal{P}_2(\cdot\,|\,\mathcal{Q}_1)) \right] \notag \\
% & \qquad
%         +\; \log\frac{m}{\delta}
%     \Big]
% \right\}^{1/2}
% \label{eq:cl-bound-chain}
% \end{align}

% The corresponding prior and hyperprior for the second task are denoted by $P_2,\mathcal{P}_2$.  In the continual learning,  the second-task hyperprior $\mathcal{P}_2$ is conditionally dependent on the hyperposterior of the first task $\mathcal{Q}_1$,i.e., $\mathcal{P}_2 = \mathcal{P}_2(\cdot|\mathcal{Q}_1)$.
% Thus, the full hyperprior for two tasks can be written as
% $\mathcal{P}_{1:2} = \mathcal{P}_1 \otimes \mathcal{P}_2(\cdot\,|\,\mathcal{Q}_1).$The corresponding hyperposterior factorizes similarly:
% $\mathcal{Q}_{1:2} = \mathcal{Q}_1 \otimes \mathcal{Q}_2.$ The joint KL divergence between the hyperdistributions takes the following chain-structured form:
% $$
% KL(\mathcal{Q}_{1:2}\,\|\,\mathcal{P}_{1:2}) = KL(\mathcal{Q}_1\,\|\,\mathcal{P}_1) + \mathbb{E}_{\mathcal{Q}_1}\left[ KL(\mathcal{Q}_2\,\|\,\mathcal{P}_2(\cdot\,|\,\mathcal{Q}_1)) \right].
% $$
% Analogously, for the parameter-level distributions, the KL divergence decomposes as
% $$KL(Q_{1:2}\,\|\,P_{1:2}) = KL(Q_1\,\|\,P_1) + KL(Q_2\,\|\,P_2),$$
% where $P_2$ is sampled from the hyperprior $\mathcal{P}_2(\cdot,|,\mathcal{Q}_1)$.










% And $\mathcal{Q}_1$ are related to the poster paramter distributon $Q_1$. Because $\mathcal{P}_2$ is conditionally dependent on $\mathcal{Q}_1$, the joint KL divergence between the hyperdistributions takes the following chain-structured form:
% $$KL(\mathcal{Q}_{1:2}\,\|\,\mathcal{P}_{1:2}) = KL(\mathcal{Q}_1\,\|\,\mathcal{P}_1) + \mathbb{E}_{\mathcal{Q}_1}\left[ KL(\mathcal{Q}_2\,\|\,\mathcal{Q}_1) \right].$$ 
% Analogously, for the parameter-level distributions, the KL divergence decomposes as
% $KL(Q_{1:2}\,\|\,P_{1:2}) = KL(Q_1\,\|\,P_1) + KL(Q_2\,\|\,P_2),$
% where $P_2$ is sampled from the hyperprior $\mathcal{P}_2$.













% % This decomposition holds because the KL divergence between independent distributions is equal to the sum of their individual KL divergences. Similarly, for hyperdistributions:
% % \begin{equation}
% %      KL(Q_{1:n}^{\text{full}} \| P_{1:n}^{\text{full}}) = \sum_{t=1}^n KL(Q_t \| P_t) .
% % \end{equation}
%  The corresponding prior, posterior, hyperprior, and hyperposterior for the second task are denoted by $P_2, Q_2(P_2), \mathcal{P}_2,$ and $\mathcal{Q}_2$, respectively. In this setting, the full posterior and prior over all parameters naturally factorize as $Q_{1:2} = Q_1 \otimes Q_2$ and $P_{1:2} = P_1 \otimes P_2$, and the same holds for the hyperdistributions. By independence, the KL divergence between the joint distributions decomposes additively:
% $$
% KL(Q_{1:2}\|P_{1:2}) = KL(Q_1\|P_1) + KL(Q_2\|P_2),
% $$
% and similarly for the hyperdistributions,
% $$
% KL(\mathcal{Q}_{1:2}\|\mathcal{P}_{1:2}) = KL(\mathcal{Q}_1\|\mathcal{P}_1) + KL(\mathcal{Q}_2\|\mathcal{P}_2).
% $$

% Thus, the generalization bound for two tasks takes the following form:
% \begin{align}
% L(\mathcal{Q}_2) :=\; &\; \mathbb{E}_{P_{1:2} \sim \mathcal{Q}_{1:2}}
%     \bigg\{
%     \frac{1}{2} \Big[
%         \mathbb{E}_{\Delta\mathbf{W}_1 \sim Q_1(P_1)}\big[
%             L_{\mathcal{D}_1}(\mathbf{W}_0 + \Delta\mathbf{W}_1)
%         \big] \notag\\
%     &\quad
%         +\;
%         \mathbb{E}_{\Delta\mathbf{W}_2 \sim Q_2(P_2)}\big[
%             L_{\mathcal{D}_2}(\mathbf{W}_1 + \Delta\mathbf{W}_2)
%         \big]
%     \Big] \bigg\}
% \label{eq:cl_exp_risk}
% \end{align}

% \begin{align}
% L(\mathcal{Q}_2) \leq &\;\; \frac{1}{2} \sum_{t=1}^2 \mathbb{E}_{P_{1:2} \sim \mathcal{Q}_{1:2}}
% \Big[
%     \widehat{L}_{S_t}(\mathbf{W}_{t-1} + \hat{\Delta\mathbf{W}}_t)
%     + \mathrm{Sharpness}_t(\hat{\Delta\mathbf{W}}_t, \rho_t^2)
% \Big] \notag \\
% & + (b-a) \sqrt{
%     \frac{
%         \sum_{t=1}^2 \mathbb{E}_{P_{1:2} \sim \mathcal{Q}_{1:2}}
%         \Big[ KL(Q_t(P_t)\|P_t) \Big]
%         + KL(\mathcal{Q}_{1:2}\|\mathcal{P}_{1:2})
%         + \log(m/\delta)
%     }{2(m-1)}
% }
% \label{eq:cl_bound}
% \end{align}


% \begin{equation}
% \begin{aligned} 
% L(\mathcal{Q}_2) \leq &\;\; \frac{1}{2}\sum_{t=1}^2 \mathbb{E}_{P_{1:2} \sim \mathcal{Q}_{1:2}}\left[\widehat{L}_{S_t}(\mathbf{W}_{t-1} + \hat{\Delta\mathbf{W}}_t) + \mathrm{Sharpness}_t(\hat{\Delta\mathbf{W}}_t, \rho_t^2)\right] \\ &+ (b-a) \sqrt{ \frac{ \sum_{t=1}^2 \mathbb{E}_{P_{1:2} \sim \mathcal{Q}_{1:2}}[KL(Q_t(P_t)\|P_t)] + KL(\mathcal{Q}_{1:2}\|\mathcal{P}_{1:2}) + \log(m/\delta) }{2(m-1)} } 
% \end{aligned}
% \end{equation}

% \begin{align}
% L(\mathcal{Q}_2) \leq\; & \frac{1}{2}\sum_{t=1}^2 
%     \mathbb{E}_{P_{1:2} \sim \mathcal{Q}_{1:2}} \Big[
%         \widehat{L}_{S_t}(\mathbf{W}_{t-1} + \hat{\Delta\mathbf{W}}_t) 
%     \notag \\
% & \qquad +\; \mathrm{Sharpness}_t(\mathbf{W}_{t-1} +\hat{\Delta\mathbf{W}}_t, \rho_t^2)
%     \Big] 
%     \notag \\
% & + (b-a) \Bigg\{
%     \frac{1}{2(\bar{m}_2-1)} \Bigg[
%         \mathbb{E}_{P_1 \sim \mathcal{Q}_1} \big[ KL(Q_1(P_1)\|P_1) \big] 
%     \notag \\
% & \qquad
%         + \mathbb{E}_{P_2 \sim \mathcal{Q}_2} \big[ KL(Q_2(P_2)\|P_2) \big] 
%     \notag \\
% & \qquad
%         + KL(\mathcal{Q}_1\|\mathcal{P}_1)
%         + KL(\mathcal{Q}_2\|\mathcal{P}_2)
%     \notag \\
% & \qquad
%         + \log\frac{m}{\delta}
%     \Bigg]
% \Bigg\}^{1/2}
% \label{eq:cl-bound-tight}
% \end{align}



% where $\bar{m}_2 = 2 / \left(\frac{1}{m_1} + \frac{1}{m_2}\right)$ denotes the harmonic mean of the sample sizes, , and $\operatorname{const}$ subsumes all terms independent of the learned parameters.

% For clarity, all terms involving empirical risk, sharpness, and complexity should be understood as being averaged over the hyperposterior, i.e., all $\hat{\Delta\mathbf{W}}t$, $\rho_t^2$ and related terms are implicitly inside the expectation over $P_{1:2} \sim \mathcal{Q}_{1:2}$. This ensures that both sides of the bound refer to the same random variable, as in the single-task case.



  
 





% \section{Appendix B: Proof of the theorem}


% \begin{lemma} \cite{pentinaPACBayesianBoundLifelong2014} \label{lemma: pentina}
%  Let $f$ be a random variable taking values in $A$ and let $X_1, \ldots, X_l$ be $l$ independent random variables with each $X_k$ distributed according to $\mu_k$ over the set $A_k$. For functions $g_k: A \times A_k \rightarrow\left[a_k, b_k\right], k=1 \ldots l$, let $\xi_k(f)=\mathbf{E}_{X_k \sim \mu_k} g_k\left(f, X_k\right)$ for any fixed value of $f$. Then for any fixed distribution $\pi$ on $A$ and any $\lambda, \delta>0$ the following inequality holds with probability at least $1-\delta$ (over sampling $X_1, \ldots, X_l$ ) for all distributions $\rho$ over $A$

% $$
% \begin{gathered}
% \underset{f \sim \rho}{\mathbf{E}} \sum_{k=1}^l \xi_k(f)-\underset{f \sim \rho}{\mathbf{E}} \sum_{k=1}^l g_k\left(h, X_k\right) \leq \\
% \frac{1}{\lambda}\left(K L(\rho \| \pi)+\frac{\lambda^2}{8} \sum_{k=1}^l\left(b_k-a_k\right)^2-\log \delta\right) .
% \end{gathered}
% $$
% \end{lemma}

% \begin{theorem} \cite{pentinaPACBayesianBoundLifelong2014}
% For any $\delta>0$ the following inequality holds with probability at least $1-\delta$ (over the training samples $\left\{S_1, \ldots, S_n\right\}$ ) for all hyperposterior distributions $\mathcal{Q}$

% $$
% \begin{aligned}
% & \operatorname{er}(\mathcal{Q}) \leq \widehat{\operatorname{er}}(\mathcal{Q})+\frac{1}{\sqrt{n}}\left(K L(\mathcal{Q} \| \mathcal{P})+\frac{1}{8}-\log \frac{\delta}{2}\right) \\
% & +\frac{1}{n \sqrt{\bar{m}}} K L\left(\left(\mathcal{Q}, Q^n\right) \|\left(\mathcal{P}, P^n\right)\right)+\frac{1}{\sqrt{\bar{m}}}\left(\frac{1}{8}-\frac{1}{n} \log \frac{\delta}{2}\right)
% \end{aligned}
% $$
% where $\left(\mathcal{Q}, Q^n\right)=\mathcal{Q} \times \prod_{i=1}^n Q_i$ denotes the distribution in which we first sample $P$ according to $\mathcal{Q}$ and then use it and the data $S_i$ to produce a posterior $Q_i$ for each task $t_i$. $\left(\mathcal{P}, P^n\right)=\mathcal{P} \times \prod_{i=1}^n P$ denotes the distribution in which we sample $P$ according to $\mathcal{P}$ and use it as a posterior for all tasks. $\bar{m}=\left(\frac{1}{n} \sum_{i=1}^n \frac{1}{m_i}\right)^{-1}$ is the harmonic mean of the sample sizes.
% \end{theorem}


% \begin{theorem} \cite{foretSharpnessAwareMinimizationEfficiently2021}
% Let $\sigma_Q=\rho$, $\mu_Q=w$, $\mu_P = 0$ and $\sigma^2_P = c \exp((1 - j)/k)$, where $c = \rho^2(1 + \exp(4m/k))$ and $j=\left\lfloor 1-k \log \left(\left(\rho^2+\|\boldsymbol{w}\|_2^2 / k\right) / c\right)\right\rfloor$. For any $\rho > 0$ and any distribution $\mathcal{D}$, with probability $1 - \delta$ over the choice of the training set $S \sim \mathcal{D}$,
% \begin{equation}
% \begin{aligned}
% &\mathbb{E}_{\epsilon_i \sim \mathcal{N}(0, \sigma)}\left[L_{\mathcal{D}}(\boldsymbol{w}+\boldsymbol{\epsilon})\right] 
% \leq 
% \mathbb{E}_{\epsilon_i \sim \mathcal{N}(0, \sigma)}\left[L_{\mathcal{S}}(\boldsymbol{w}+\boldsymbol{\epsilon})\right] \\
% &+\sqrt{\frac{\frac{1}{4} k \log \left(1+\frac{\|\boldsymbol{w}\|_2^2}{k \sigma^2}\right)+\frac{1}{4}+\log \frac{m}{\delta}+2 \log (6 m+3 k)}{m-1}}
% \end{aligned}
% \end{equation}
    
% \end{theorem}




% \begin{theorem} 
% For any confidence parameter  $ \delta > 0  $, with probability at least  $ 1 - \delta  $ over the sampling of training sets  $ \{ S_1, \dots, S_n \}  $, the following PAC-Bayesian generalization bound holds for all hyperposterior distributions  $ \mathcal{Q}  $: 
% \[ \begin{aligned} er(\mathcal{Q}) &\leq \frac{1}{n} \sum_{i=1}^n \hat{L}_{S_i}(Q_i) + \frac{1}{n} \sum_{i=1}^n \text{Sharpness}(Q_i) \\ &+ \frac{1}{\sqrt{n}} \left( KL(\mathcal{Q} \| \mathcal{P}) + \frac{1}{8} - \log \frac{\delta}{2} \right) \\ &+ \frac{1}{n\sqrt{\bar{m}}} KL\left( (\mathcal{Q}, Q^n) \| (\mathcal{P}, P^n) \right) + \frac{1}{\sqrt{\bar{m}}} \left( \frac{1}{8} - \frac{1}{n} \log \frac{\delta}{2} \right), \end{aligned} \] 
% \end{theorem}


% \begin{proof}

% The proof technique we use here is inspired from \cite{pentinaPACBayesianBoundLifelong2014} and \cite{foretSharpnessAwareMinimizationEfficiently2021}. We also use the lemma from pentina \shortcite{pentinaPACBayesianBoundLifelong2014} and conclusion from foret \shortcite{foretSharpnessAwareMinimizationEfficiently2021} in as above. 

% To prove Theorem 1 we introduce the continual generalization risk with perturbation
% \begin{equation}
%     L(\mathcal{Q}) = \mathbb{E}_{(t,S_t)} \mathbb{E}_{P \sim \mathcal{Q}} \mathbb{E}_{f \sim Q_t} \mathbb{E}_{\epsilon \sim \mathcal{N}(0, \sum_{Q_t})} \left[ L_{\mathcal{D}_t}( f+ \epsilon) \right],
% \end{equation}
% the average generalization risk on observed tasks with Perturbation
% \begin{equation}
% \begin{aligned}
%     L(\mathcal{Q}_{1:n}) :=\; & \frac{1}{n} \sum_{t=1}^n  
%     \mathbb{E}_{P_{1:n} \sim \mathcal{Q}_{1:n}}\,
%     \mathbb{E}_{f \sim Q_t(P_t)}\, \\
%     &\qquad
%     \mathbb{E}_{\epsilon \sim \mathcal{N}(0,\,\Sigma_{Q_t})} 
%     \left[ L_{\mathcal{D}_t}(f + \epsilon) \right]
% \end{aligned}
% \end{equation}
% and  the Continual Empirical Risk with Perturbation
% \begin{equation}
% \begin{aligned}
%     \widehat{L}(\mathcal{Q}_{1:n}) :=\; & \frac{1}{n} \sum_{t=1}^n
%     \mathbb{E}_{P \sim \mathcal{Q}}\, \mathbb{E}_{f \sim Q_t}\, \\
%     &\quad
%     \mathbb{E}_{\epsilon \sim \mathcal{N}(0,\,\Sigma_{Q_t})}
%     \left[ \frac{1}{m_t} \sum_{j=1}^{m_t} \ell(f(x_{tj}) + \epsilon,\, y_{tj}) \right]
% \end{aligned}
% \end{equation}

% In order to bound the difference between $L_{\epsilon}(\mathcal{Q})$ and $\tilde{L}_\epsilon(\mathcal{Q})$ we treat each task $t$ with the corresponding training sample $S_t$ as a random variable and apply Lemma \ref{lemma: pentina} from pentina \shortcite{pentinaPACBayesianBoundLifelong2014}. Formally, we set $f = \sigma^2$, $\rho = \mathcal{Q}$, $\pi = \mathcal{P}$, $X_k = (t_k, S_k)$, $l = n$, $ g_k(f, X_k) = \frac{1}{n} \mathbb{E}_{h \sim Q_k} \mathbb{E}_{\epsilon \sim \mathcal{N}(0, f\mathbf{I})} [\ell(h(x) + \epsilon, y)] $, and using $\lambda = \sqrt{n}$, we obtain (with probability at least $1-\delta/2$):
% $$  L_{\epsilon}(\mathcal{Q}) \leq  L_\epsilon(\mathcal{Q}_{1:n}) + \frac{1}{\sqrt{n}}\left( \mathrm{KL}(\mathcal{Q} \| \mathcal{P}) + \frac{1}{8} - \log\frac{\delta}{2} \right).$$





% \end{proof}





\end{document}
